\section{核心发现}

\begin{reportbox}{分析摘要}
本报告围绕近五年公开科技文献展开结构化分析。当前样本覆盖 6 个公开来源、169,324 条原始记录和 168,447 条分析语料, 其中 159,382 条位于 2022--2026 年主分析窗口。整体上, 年份覆盖较均衡, 领域结构以计算机科学和生物医学为主, 材料科学样本相对较少; 因此主题变化和领域比较均采用归一化口径辅助解释。
\end{reportbox}

\begin{center}
\begin{tabularx}{\textwidth}{>{\bfseries}p{34mm}X}
\toprule
分析问题 & 主要发现 \\
\midrule
数据规模与覆盖 & 分析语料达到 168,447 条, 2022--2026 每年均超过 30,000 条分析记录, 可支撑近五年趋势观察。 \\
领域结构 & 计算机科学 75,493 条、生物医学 64,858 条、材料科学 28,096 条; 材料科学占比较低, 需要避免只用绝对频次判断热度。 \\
热点主题 & retrieval augmented generation、large language model、prompt engineering、vision language model、knowledge graph 等方向在融合术语、归一化趋势、共现结构和主题模型中同时出现。 \\
作者合作 & 作者合作网络包含 10,499,041 条候选合作边; 分数计数后仍能识别高合作度作者、桥接作者和稳定社区。 \\
质量边界 & 摘要覆盖率为 83.22\%, 正文片段覆盖率为 8.99\%; 摘要级结论可信度较高, 全文级深读仍受来源差异影响。 \\
\bottomrule
\end{tabularx}
\end{center}

\subsection{可直接采用的结论}

\begin{itemize}
  \item \textbf{规模结论}: 当前数据已经满足近五年、多来源、多领域的统计分析需要。
  \item \textbf{趋势结论}: LLM、RAG、prompt engineering、vision language model 等主题具有明确上升信号; 这些结论来自显式关键词、题名摘要术语、共现网络和主题模型的交叉支持, 但仍应表述为候选趋势而非确定预测。
  \item \textbf{网络结论}: 作者合作网络显示出明显社区结构, 适合用于识别合作团簇和跨领域连接者。
  \item \textbf{质量结论}: 来源偏差、关键词噪声和作者同名问题会影响细粒度排名, 因此核心结论应以趋势和结构为主, 少做单点绝对判断。
  \item \textbf{价值结论}: 报告可服务科研管理、科技情报、产业研发、投资初筛和政策制定, 将分散文献记录转化为可复核的热点监测和方向研判证据。
  \item \textbf{创新结论}: 报告采用融合关键词信号、年度归一化、动量/burst、生命周期、共现网络、主题模型和质量审计的多证据框架, 避免只凭单一词频榜判断热点。
\end{itemize}

\clearpage
